\documentclass[oneside]{utmsrs}
% Options
% Size [a4draft, b5final]
% Language [english, malay]
% Printing [twoside, oneside]
% Default option [a4draft, english, twoside]

\usepackage{graphicx}
\usepackage{url} 
\usepackage{pdflscape}
\usepackage{longtable}
\usepackage{notoccite}
\usepackage[numbers]{natbib}
\usepackage{multirow}%
\usepackage{algorithm}
\usepackage{algpseudocode}
\usepackage{array}%
\usepackage{color,soul}%
\usepackage[table,xcdraw]{xcolor}
\usepackage{booktabs}
\usepackage[caption=false]{subfig}
\newcolumntype{P}[1]{>{\centering\arraybackslash}m{#1}}%

%This is to make sure vertical spacing non-stretchable
\raggedbottom

% Required information
% \title{Software Requirements Specification}
% \subtitle{SECJ 3032: Software Engineering FYP1}
% \faculty{FACULTY OF COMPUTING}
% \university{UNIVERSITI TEKNOLOGI MALAYSIA}
% \location{UTM Johor Bahru}
% \semester{Semester 01, 2024/2025}

\begin{document}

% Custom front page
\begin{titlepage}
    \begin{center}
        \vspace*{1cm}
        
        \textbf{\Large UTM}
        
        \vspace{0.5cm}
        
        \textbf{\large FACULTY OF COMPUTING}
        
        \vspace{0.5cm}
        
        \textbf{\large UTM Johor Bahru}
        
        \vspace{1cm}
        
        \textbf{\large SECJ 3032: Software Engineering FYP1}
        
        \vspace{0.5cm}
        
        \textbf{\large Semester 01, 2024/2025}
        
        \vspace{1.5cm}
        
        \textbf{\Huge Software Requirements Specification}
        
        \vspace{0.5cm}
        
        \textbf{\Large (SRS)}
        
        \vspace{2cm}
        
        \textbf{\LARGE <<Project Title>>}
        
        \vspace{1cm}
        
        \textbf{\large <<Version No.>>}
        
        \vspace{2cm}
        
        \textbf{\large Prepared by: <<Name>>}
        
        \vspace*{\fill}
        
        \textbf{\large \today}
    \end{center}
\end{titlepage}

% Rest of the document remains the same
\frontmatter

% Revision Page
\chapter*{Revision Page}
\section*{Overview}
This document specifies the software requirements for <<Project Title>>. It describes the functional and non-functional requirements, system features, and design constraints.

\section*{Target Audience}
The intended audience for this document includes:
\begin{itemize}
    \item Project stakeholders
    \item Software developers
    \item Testers
    \item Project supervisors
\end{itemize}

\section*{Version Control History}
\begin{longtable}{|l|l|l|l|}
\hline
\textbf{Version} & \textbf{Primary Author(s)} & \textbf{Description of Version} & \textbf{Date Completed} \\ \hline
1.0 & Team leader 1 (Full Name) & Completed Chapter X, Section... & dd/mm/yyyy \\ \hline
... & ... & ... & ... \\ \hline
\end{longtable}

\textbf{Note:} This Software Requirements Specification (SRS) template is adapted from IEEE Recommended Practice for Software Requirements Specification (SRS) (IEEE Std. 830-1998) that is simplified and customized to meet the need of SECJ3203 FYP1 SE course at Faculty of Computing, UTM.

% Table of Contents
\tableofcontents

\mainmatter

\chapter{Introduction}
\section{Purpose}
This Software Requirements Specification (SRS) document:
\begin{itemize}
    \item Delineates the purpose of the SRS
    \item Specifies the intended audience for the SRS
    \item Describes the software product to be developed
\end{itemize}

\section{Scope}
The software product is <<Project Title>> which will:
\begin{itemize}
    \item Identify the software product(s) to be produced
    \item Explain what the software product(s) will and will not do
    \item Describe the application of the software being specified
    \item Be consistent with higher-level specifications
\end{itemize}

\section{Definitions, Acronyms and Abbreviations}
Definitions of all terms, acronyms and abbreviations used are to be defined here.

\section{References}
\begin{itemize}
    \item IEEE Std. 830-1998, IEEE Recommended Practice for Software Requirements Specifications
    \item Arlow and Neustadt (2002), UML and the Unified Process
\end{itemize}

\section{Overview}
This document contains:
\begin{itemize}
    \item Introduction to the project
    \item Specific requirements including user characteristics and system features
    \item Performance and design constraints
    \item Appendices with supporting materials
\end{itemize}

\chapter{Specific Requirements}
This section contains all software requirements to enable designers to design the system and testers to verify it.

\section{User Characteristics}
The system will interact with the following user types:

\subsection{<<User Type 1>>}
\begin{itemize}
    \item Expected technical skills
    \item Knowledge level
    \item Physical abilities
    \item Specific requirements
\end{itemize}

\subsection{<<User Type 2>>}
\begin{itemize}
    \item Expected technical skills
    \item Knowledge level
    \item Physical abilities
    \item Specific requirements
\end{itemize}

\section{System Features}
The system will provide the following features:
\begin{itemize}
    \item Feature 1 with description
    \item Feature 2 with description
    \item Feature n with description
\end{itemize}

% \begin{figure}[h]
%     \centering
%     \includegraphics[width=0.8\textwidth]{use_case_diagram.png}
%     \caption{Use Case Diagram for <<System Name>>}
%     \label{fig:use_case}
% \end{figure}

\section{Use Case Details}
\subsection{UC001: <<Use Case Name>>}
\begin{table}[h]
\centering
\begin{tabular}{|p{0.3\textwidth}|p{0.6\textwidth}|}
\hline
\textbf{Use Case ID} & UC001 \\ \hline
\textbf{Use Case Name} & <<Action Verb>> \\ \hline
\textbf{Description} & <<Brief description>> \\ \hline
\textbf{Actor(s)} & <<List of actors>> \\ \hline
\textbf{Pre-condition(s)} & <<List of constraints>> \\ \hline
\textbf{Normal Flow(s)} & 
1. The user performs action \\
2. The system responds \\ \hline
\textbf{Alternative Flow(s)} & AF1. Alternative scenario \\ \hline
\textbf{Exception Flow(s)} & EF1. Error handling \\ \hline
\textbf{Post-condition(s)} & <<System states>> \\ \hline
\end{tabular}
\caption{Use Case Description for <<Use Case Name>>}
\label{tab:uc001}
\end{table}

% \begin{figure}[h]
%     \centering
%     \includegraphics[width=0.8\textwidth]{sequence_diagram_uc001.png}
%     \caption{Sequence Diagram for UC001}
%     \label{fig:seq_uc001}
% \end{figure}

% \begin{figure}[h]
%     \centering
%     \includegraphics[width=0.8\textwidth]{activity_diagram_uc001.png}
%     \caption{Activity Diagram for UC001}
%     \label{fig:act_uc001}
% \end{figure}

\section{Performance and Other Requirements}
\subsection{Software System Attributes}
\begin{itemize}
    \item Usability: The system should be easy to use
    \item Reliability: The system should perform consistently
    \item Maintainability: The system should be easy to modify
    \item Portability: The system should run on multiple platforms
    \item Compatibility: The system should work with other systems
\end{itemize}

\subsection{Performance Requirements}
\begin{itemize}
    \item Response Time: Maximum 5 seconds
    \item Throughput: 1000 requests per minute
    \item Capacity: 500 concurrent users
    \item Availability: 99.9\% uptime
\end{itemize}

\subsection{Other Requirements}
\begin{itemize}
    \item Security: Data encryption required
    \item Safety: No harm to users
    \item Legal: Compliance with regulations
    \item Environmental: Low power consumption
\end{itemize}

\section{Design Constraints}
\begin{itemize}
    \item Environmental: Operating temperature range
    \item Hardware: Minimum system requirements
    \item Security: Encryption standards
    \item Compatibility: Supported platforms
    \item Performance: Maximum response times
\end{itemize}

\appendix
\chapter{Evidence of Requirements Elicitation Artefacts}
Appendix content goes here.

\end{document}